% !TeX program = pdflatex
%==============================================================================
%  A DETAILED FORMALISATION BLUEPRINT
%  From current Mathlib to Grothendieck's Coherence Theorem (EGA III.3.2.1)
%
%  Target (from derivedpushforward.lean):
%    {X Y : Scheme.{u}} (f : X ⟶ Y) [IsProper f] [IsLocallyNoetherian Y]
%    (M : X.Modules) [M.IsFinitePresentation] (n : ℕ) :
%    (((Modules.pushforward f).rightDerived n).obj M).IsFinitePresentation
%
%  This blueprint decomposes the whole development into a dependency graph of
%  definitions and lemmas, each with a statement, references, dependencies
%  (\uses), an intended Lean name (\lean), a Mathlib-status flag, and a proof
%  sketch, starting from what is in Mathlib today and ending at the theorem.
%
%  PORTING TO leanblueprint: move the body into blueprint/src/content.tex and
%  delete the STANDALONE SHIM block; the package supplies the macros and the
%  theorem environments.
%
%  Mathlib status flags:  \mathlibok  (present today, modulo API glue)
%                         \mathlibpartial (partly present; must be extended)
%                         \notready   (absent; must be built)
%  Reference tags "NNNN" are Stacks Project tags (stacks.math.columbia.edu).
%==============================================================================
\documentclass[11pt]{report}
\usepackage[margin=1in]{geometry}
\usepackage{amsmath,amssymb,amsthm}
\usepackage{mathtools}
\usepackage{alltt}
\usepackage[T1]{fontenc}
\usepackage{enumitem}
\usepackage{xcolor}
\usepackage{hyperref}
\hypersetup{colorlinks=true,linkcolor=blue,citecolor=blue,urlcolor=blue}

%--------------------- STANDALONE SHIM (delete under leanblueprint) ------------
\providecommand{\lean}[1]{\par\smallskip\noindent{\footnotesize\textsf{Lean:}\ \texttt{#1}}}
\providecommand{\uses}[1]{\par\noindent{\footnotesize\textsf{Uses:}\ \texttt{#1}}}
\providecommand{\leanok}{}
\providecommand{\notready}{\par\noindent{\footnotesize\textcolor{gray}{[absent from Mathlib: to be built]}}}
\providecommand{\mathlibok}{\par\noindent{\footnotesize\textcolor{gray}{[in Mathlib today, modulo API glue]}}}
\providecommand{\mathlibpartial}{\par\noindent{\footnotesize\textcolor{gray}{[partly in Mathlib: must be extended]}}}
\providecommand{\discussion}[1]{}
%------------------------------------------------------------------------------

\theoremstyle{plain}
\newtheorem{theorem}{Theorem}[chapter]
\newtheorem{proposition}[theorem]{Proposition}
\newtheorem{lemma}[theorem]{Lemma}
\newtheorem{corollary}[theorem]{Corollary}
\theoremstyle{definition}
\newtheorem{definition}[theorem]{Definition}
\newtheorem{construction}[theorem]{Construction}
\newtheorem{remark}[theorem]{Remark}

\newcommand{\Spec}{\operatorname{Spec}}
\newcommand{\Proj}{\operatorname{Proj}}
\newcommand{\Supp}{\operatorname{Supp}}
\newcommand{\Hom}{\operatorname{Hom}}
\newcommand{\Gammaf}{\Gamma}
\newcommand{\Ff}{\mathcal{F}}
\newcommand{\Gg}{\mathcal{G}}
\newcommand{\Hh}{\mathcal{H}}
\newcommand{\Ii}{\mathcal{I}}
\newcommand{\Ll}{\mathcal{L}}
\newcommand{\Oo}{\mathcal{O}}
\newcommand{\Bb}{\mathcal{B}}
\newcommand{\PP}{\mathbf{P}}
\newcommand{\kk}{\kappa}
\newcommand{\Rr}{\mathrm{R}}
\newcommand{\Cech}{\check{\mathrm{C}}}
\newcommand{\stacks}[1]{\href{https://stacks.math.columbia.edu/tag/#1}{\texttt{#1}}}

\title{A Detailed Formalisation Blueprint\\
\large From current Mathlib to Grothendieck's Coherence Theorem\\
(EGA III, Th\'eor\`eme 3.2.1)}
\author{Blueprint for a Lean~4 / Mathlib formalisation}
\date{\today}

\begin{document}
\maketitle
\tableofcontents

%##############################################################################
\part{Baseline, strategy, and the target}
%##############################################################################

\chapter{The target theorem}\label{chap:target}

\begin{theorem}[Grothendieck's Coherence Theorem; EGA III, Th\'eor\`eme 3.2.1]\label{thm:goal}
Let $Y$ be a locally Noetherian scheme and $f\colon X\to Y$ a proper morphism.
For every coherent $\Oo_X$-module $\Ff$ and every $n\ge 0$, the $\Oo_Y$-module
$\Rr^n f_*\Ff$ is coherent. Equivalently, in the form of
\texttt{derivedpushforward.lean}: for $f$ proper, $Y$ locally Noetherian and
$M$ a sheaf of modules of finite presentation on $X$, the object
$\bigl((\mathrm{Modules.pushforward}\,f).\mathrm{rightDerived}\,n\bigr)(M)$ is
of finite presentation for every $n$.
\lean{AlgebraicGeometry.Scheme.rightDerivedPushforward\_isFinitePresentation\_of\_isProper}
\uses{thm:finiteness, prop:bridge-rderived, lem:fp-iff-coh}
\notready
\end{theorem}

Reference: EGA III \cite{EGA3}; Stacks, Proposition 30.19.1, tag~\stacks{02O5}
\cite{stacks}; Hartshorne, III.8.8 (projective case) \cite{Hartshorne};
Vakil, \S18.9 \cite{Vakil}. The whole blueprint is the proof of
Theorem~\ref{thm:goal}: Part~\ref{part:qcoh}--\ref{part:cech} build the
missing cohomological substrate on top of what Mathlib has today,
Parts~\ref{part:proj}--\ref{part:chow} build the geometric inputs, and
Part~\ref{part:finiteness} assembles the theorem.

\chapter{What Mathlib provides today}\label{chap:baseline}

The following are present in Mathlib and are the starting line. Item names are
the relevant modules or declarations; ``modulo API glue'' means the object
exists but auxiliary lemmas about it may still be needed.

\begin{description}[leftmargin=1.4em]
\item[Schemes.] \texttt{AlgebraicGeometry.Scheme}, morphisms, \texttt{Spec},
the $\Gamma\!\dashv\!\Spec$ adjunction, fibre products, open/closed
immersions, and a large morphism-property library (finite type, separated,
quasi-compact, quasi-separated, affine, universally closed, integral, closed
immersion). \mathlibok
\item[Proper and Noetherian.] \texttt{AlgebraicGeometry.Morphisms.Proper}
(\texttt{IsProper}); \texttt{AlgebraicGeometry.Noetherian}
(\texttt{IsNoetherian}, \texttt{IsLocallyNoetherian}). \mathlibok
\item[Sheaves of modules.]
\texttt{Mathlib.Algebra.Category.ModuleCat.Sheaf} and
\texttt{AlgebraicGeometry.Modules.Sheaf}: the category
$\texttt{SheafOfModules}$, its abelian structure
(\texttt{Modules.instAbelian}), the pushforward and pullback functors
(\texttt{Scheme.Modules.pushforward}, \texttt{pullback}), the adjunction
(\texttt{pullbackPushforwardAdjunction}), and the coherence isomorphisms
\texttt{pushforwardId}, \texttt{pushforwardComp}. \mathlibok
\item[Quasi-coherence, abstractly.]
\texttt{Mathlib.Algebra.Category.ModuleCat.Sheaf.Quasicoherent}: a sheaf of
modules is quasi-coherent if it is locally the cokernel of a map between
coproducts of copies of the structure sheaf, and of finite presentation when
these are finite (\texttt{SheafOfModules.IsFinitePresentation}). \mathlibok
\item[The affine model $\widetilde M$.]
\texttt{AlgebraicGeometry.Modules.Tilde}: for an $R$-module $M$, the sheaf
$\widetilde M$ of $\Oo_{\Spec R}$-modules, with its localisation-to-stalk
maps. \mathlibpartial\ (the sheaf and its stalks exist; the equivalence
$\mathrm{QCoh}(\Spec R)\simeq \mathrm{Mod}_R$ and the pushforward/pullback
compatibilities are not all present).
\item[Derived functors.]
\texttt{CategoryTheory.Abelian.RightDerived} and
\texttt{Abelian.GrothendieckCategory.EnoughInjectives}: any Grothendieck
abelian category (in particular \texttt{SheafOfModules}) has enough
injectives, so \texttt{Functor.rightDerived} of a left exact functor is
defined; hence
$(\mathrm{Modules.pushforward}\,f).\mathrm{rightDerived}\,n$ exists. \mathlibok
\item[Proj.] \texttt{AlgebraicGeometry.ProjectiveSpectrum.Scheme}: $\Proj S$
of a graded ring as a scheme. \mathlibpartial\ (the scheme exists; the
twisting sheaves and the graded-module$\,\to\,$sheaf functor are absent).
\item[Homological algebra.] Abelian categories, injective resolutions,
$\delta$-functors, long exact sequences, $\mathrm{Ext}$, spectral objects and
some spectral-sequence API (\texttt{CategoryTheory.Abelian},
\texttt{Homology}). \mathlibok
\end{description}

\chapter{What must be built, and the critical path}\label{chap:gaps}

Relative to Chapter~\ref{chap:baseline}, the development requires the
following, in dependency order. Each is a Part below.
\begin{enumerate}[leftmargin=1.6em]
\item \textbf{Quasi-coherent sheaves on schemes} (Part~\ref{part:qcoh}): the
affine equivalence via $\widetilde{(-)}$, quasi-coherence of the geometric
$\Oo_X$-modules, and the operations used later.
\item \textbf{Cohomology of $\Oo_X$-modules} (Part~\ref{part:cohomology}):
$H^i(X,-)$ and $\Rr^i f_*$ as derived functors, the flasque package, and the
comparison of the abstract \texttt{rightDerived} with geometric cohomology.
\item \textbf{\v{C}ech cohomology and affine vanishing} (Part~\ref{part:cech}):
the \v{C}ech complex, Serre's vanishing on affines, \v{C}ech computes
cohomology for affine covers, and quasi-coherence of $\Rr^i f_*$.
\item \textbf{Projective space and Proj} (Part~\ref{part:proj}): twisting
sheaves, coherent sheaves from graded modules, and the explicit cohomology of
$\PP^n$.
\item \textbf{Coherent sheaves} (Part~\ref{part:coherent}): the coherence
predicate, coherence $=$ finite presentation on locally Noetherian schemes,
operations, closed immersions, support.
\item \textbf{Projective morphisms} (Part~\ref{part:projmor}): relative
ampleness and coherence/vanishing of $\Rr^p f_*$ along projective morphisms.
\item \textbf{D\'evissage} (Part~\ref{part:devissage}) and \textbf{Chow's
lemma} (Part~\ref{part:chow}).
\item \textbf{The finiteness theorem} (Part~\ref{part:finiteness}) and the
\textbf{bridge} to the Lean statement (Part~\ref{part:bridge}).
\end{enumerate}

\begin{remark}[On spectral sequences]\label{rem:no-ss}
The relative Leray (Grothendieck) spectral sequence is \emph{not} on the
critical path. Following the note \cite{Note} incorporated earlier, the single
step of the [EGA] proof that used it is carried out by the flasque
acyclic-pushforward Lemma~\ref{lem:acyclic-pushforward}, which lives in the
flasque layer of Part~\ref{part:cohomology} and needs only: injectives are
flasque, pushforward preserves flasqueness, flasque sheaves are acyclic, and
acyclic resolutions compute derived functors. No construction of a spectral
sequence of a composite of functors is required anywhere below.
\end{remark}

\chapter{Conventions}\label{chap:conventions}
Schemes and morphisms are as in Mathlib's \texttt{AlgebraicGeometry}. For an
$\Oo_X$-module $\Ff$, $H^i(X,\Ff)$ is the $i$-th right derived functor of
$\Gamma(X,-)$ on the abelian category of $\Oo_X$-modules, and $\Rr^i f_*$ the
$i$-th right derived functor of $f_*$. ``Coherent'' is
Definition~\ref{def:coherent}; on locally Noetherian schemes it coincides with
quasi-coherent of finite type and with finite presentation
(Lemma~\ref{lem:fp-iff-coh}). A \emph{projective morphism} is one admitting a
closed immersion into some $\PP^n_S$ over the base
(Definition~\ref{def:projective-morphism}).

%##############################################################################
\part{Quasi-coherent sheaves on schemes}\label{part:qcoh}
%##############################################################################

\chapter{The affine model: the functor \texorpdfstring{$\widetilde{(-)}$}{tilde}}\label{chap:tilde}

Reference: Stacks, Section 26.7, tag~\stacks{01I6}, and Section 17.10
\cite{stacks}; Hartshorne, II.5 \cite{Hartshorne}.

\begin{construction}[The functor $\widetilde{(-)}$]\label{con:tilde}
For a commutative ring $R$, sending an $R$-module $M$ to the
$\Oo_{\Spec R}$-module $\widetilde M$ (Mathlib's
\texttt{ModuleCat.tildeInModuleCat}) is functorial and additive, giving
$\widetilde{(-)}\colon \mathrm{Mod}_R\to \mathrm{Mod}(\Oo_{\Spec R})$.
\lean{ModuleCat.tildeFunctor}
\mathlibpartial
\end{construction}

\begin{lemma}[Sections and stalks of $\widetilde M$]\label{lem:tilde-sections}
For $f\in R$ and a prime $\mathfrak p$,
$\widetilde M(D(f))\cong M_f$ (as $R_f$-modules), $\Gamma(\Spec R,\widetilde
M)\cong M$, and $(\widetilde M)_{\mathfrak p}\cong M_{\mathfrak p}$; these are
natural in $M$. Reference: Stacks, Lemma 26.7.1, tag~\stacks{01I6}.
\lean{ModuleCat.Tilde.sections\_basicOpen, ModuleCat.Tilde.stalkIso}
\uses{con:tilde}
\mathlibpartial
\end{lemma}

\begin{lemma}[$\widetilde{(-)}$ is exact and fully faithful]\label{lem:tilde-ff}
The functor $\widetilde{(-)}$ is exact and fully faithful; on morphisms it is
the bijection $\Hom_R(M,N)\cong\Hom_{\Oo_{\Spec R}}(\widetilde M,\widetilde
N)$. Reference: Stacks, tag~\stacks{01I6}; Hartshorne, II.5.2, II.5.5.
\lean{ModuleCat.tildeFunctor\_fullyFaithful, ModuleCat.tildeFunctor\_exact}
\uses{lem:tilde-sections}
\notready
\end{lemma}

\begin{proposition}[Affine equivalence]\label{prop:affine-equiv}
$\widetilde{(-)}$ induces an equivalence
$\mathrm{Mod}_R\simeq \mathrm{QCoh}(\Spec R)$ onto the quasi-coherent
$\Oo_{\Spec R}$-modules, with quasi-inverse $\Gamma(\Spec R,-)$. Reference:
Stacks, Lemma 26.7.1 and Section 26.7, tag~\stacks{01I6}.
\lean{AlgebraicGeometry.Spec.qcohEquivModule}
\uses{lem:tilde-ff, def:qcoh}
\notready
\end{proposition}

\begin{lemma}[Functoriality on affines]\label{lem:tilde-functorial}
For a ring map $\varphi\colon R\to S$ inducing $\psi\colon \Spec S\to\Spec R$:
$\psi_*\widetilde N\cong\widetilde{N_R}$ (restriction of scalars) for an
$S$-module $N$, and $\psi^*\widetilde M\cong\widetilde{S\otimes_R M}$ for an
$R$-module $M$, naturally. Reference: Stacks, Lemma 26.7.3, tag~\stacks{01I6}.
\lean{AlgebraicGeometry.Spec.pushforward\_tilde, AlgebraicGeometry.Spec.pullback\_tilde}
\uses{prop:affine-equiv}
\notready
\end{lemma}

\chapter{Quasi-coherent sheaves on general schemes}\label{chap:qcoh}

\begin{definition}[Quasi-coherent $\Oo_X$-module]\label{def:qcoh}
An $\Oo_X$-module $\Ff$ is \emph{quasi-coherent} if $X$ has an affine open
cover $\{U_i=\Spec R_i\}$ with $\Ff|_{U_i}\cong\widetilde{M_i}$ for
$R_i$-modules $M_i$. Reference: Stacks, tag~\stacks{01BD}; Hartshorne, II.5.4.
\lean{AlgebraicGeometry.Scheme.Modules.IsQuasicoherent}
\uses{con:tilde}
\notready
\end{definition}

\begin{lemma}[Agreement with the abstract predicate]\label{lem:qcoh-agree}
Definition~\ref{def:qcoh} agrees with the abstract quasi-coherence of
\texttt{SheafOfModules.IsQuasicoherent} (local presentation as a cokernel of
free modules): a geometric $\Oo_X$-module is quasi-coherent in the sense of
Definition~\ref{def:qcoh} iff its image in \texttt{SheafOfModules} is. This is
the link between the geometric layer and the existing Mathlib predicate.
Reference: Stacks, tag~\stacks{01BG}.
\lean{AlgebraicGeometry.Scheme.Modules.isQuasicoherent\_iff\_sheafOfModules}
\uses{def:qcoh, prop:affine-equiv}
\notready
\end{lemma}

\begin{lemma}[Local nature]\label{lem:qcoh-local}
Quasi-coherence is local: if $\Ff|_{U_i}$ is quasi-coherent for an open cover
$\{U_i\}$ then $\Ff$ is; and on an affine $\Spec R$ every quasi-coherent sheaf
is $\widetilde M$ for $M=\Gamma$. Reference: Stacks, tag~\stacks{01I6}.
\lean{AlgebraicGeometry.Scheme.Modules.isQuasicoherent\_of\_cover}
\uses{def:qcoh, prop:affine-equiv}
\notready
\end{lemma}

\begin{lemma}[The structure sheaf and pullbacks]\label{lem:qcoh-basic}
$\Oo_X$ is quasi-coherent, and for any $f\colon X\to Y$ the pullback
$f^*\Gg$ of a quasi-coherent $\Gg$ is quasi-coherent. Reference: Stacks,
tags~\stacks{01BE}, \stacks{01BM}.
\lean{AlgebraicGeometry.Scheme.Modules.isQuasicoherent\_pullback}
\uses{def:qcoh, lem:tilde-functorial}
\notready
\end{lemma}

\begin{proposition}[Quasi-coherent sheaves form a Serre--abelian subcategory]\label{prop:qcoh-abelian}
Kernels, cokernels, images and finite direct sums of morphisms of
quasi-coherent sheaves are quasi-coherent, as are extensions; hence
$\mathrm{QCoh}(X)$ is a weak Serre subcategory of $\mathrm{Mod}(\Oo_X)$ and is
abelian, and the inclusion is exact. Reference: Stacks, tags~\stacks{01BJ},
\stacks{01LA}.
\lean{AlgebraicGeometry.Scheme.Modules.qcoh\_isAbelian}
\uses{def:qcoh, lem:tilde-ff, prop:affine-equiv}
\notready
\end{proposition}

\begin{remark}
Quasi-coherence of \emph{pushforwards} $f_*$ and of $\Rr^i f_*$ needs
cohomological input (vanishing on affines) and is deferred to
Chapter~\ref{chap:hdi-qcoh}.
\end{remark}

\chapter{Invertible sheaves and twists}\label{chap:invertible}

Only a light amount is needed, for the projective case and Chow's lemma.

\begin{definition}[Locally free and invertible]\label{def:invertible}
An $\Oo_X$-module is \emph{locally free of rank $r$} if $X$ has an open cover
on which it is $\Oo^{\oplus r}$, and \emph{invertible} if locally free of rank
$1$. Invertible sheaves form a group $\mathrm{Pic}(X)$ under $\otimes_{\Oo_X}$.
Reference: Stacks, tag~\stacks{01CR}; Hartshorne, II.6.12.
\lean{AlgebraicGeometry.Scheme.Modules.IsInvertible}
\mathlibpartial
\end{definition}

\begin{lemma}[Tensoring by an invertible sheaf is exact]\label{lem:tensor-invertible-exact}
For an invertible $\Ll$, the functor $-\otimes_{\Oo_X}\Ll$ is exact and
preserves quasi-coherence and (later) coherence, with inverse
$-\otimes\Ll^{-1}$. Reference: Stacks, tag~\stacks{01CR}.
\lean{AlgebraicGeometry.Scheme.Modules.tensorInvertible\_exact}
\uses{def:invertible, prop:qcoh-abelian}
\notready
\end{lemma}

%##############################################################################
\part{Cohomology of \texorpdfstring{$\Oo_X$}{OX}-modules}\label{part:cohomology}
%##############################################################################

References: Stacks, Chapter 20 (Cohomology of Sheaves) and Chapter 30 (Cohomology
of Schemes) \cite{stacks}; Hartshorne, III.1--III.2, III.8 \cite{Hartshorne}.

\chapter{Sheaf cohomology and higher direct images}\label{chap:cohomology-def}

\begin{definition}[Sheaf cohomology]\label{def:sheaf-cohomology}
For a scheme $X$, $H^i(X,-)$ is the $i$-th right derived functor of
$\Gamma(X,-)\colon \mathrm{Mod}(\Oo_X)\to \mathrm{Ab}$, computed by injective
resolutions (which exist: $\mathrm{Mod}(\Oo_X)$ is a Grothendieck abelian
category, already available in Mathlib). Reference: Stacks, tag~\stacks{01FR};
Hartshorne, III.2.
\lean{AlgebraicGeometry.Scheme.cohomology}
\uses{}
\notready
\end{definition}

\begin{definition}[Higher direct image]\label{def:higher-direct-image}
For $f\colon X\to Y$, $\Rr^i f_*$ is the $i$-th right derived functor of the
left exact functor $f_*\colon \mathrm{Mod}(\Oo_X)\to\mathrm{Mod}(\Oo_Y)$.
Reference: Stacks, tag~\stacks{01F1}; Hartshorne, III.8.
\lean{AlgebraicGeometry.Scheme.higherDirectImage}
\uses{}
\notready
\end{definition}

\begin{lemma}[Left exactness]\label{lem:pf-left-exact}
$\Gamma(X,-)$ and $f_*$ are additive and left exact; $f_*$ is right adjoint to
$f^*$ (present in Mathlib as \texttt{pullbackPushforwardAdjunction}), whence
left exactness. Reference: Stacks, tag~\stacks{01F1}.
\lean{AlgebraicGeometry.Scheme.Modules.pushforward\_leftExact}
\uses{def:higher-direct-image}
\mathlibpartial
\end{lemma}

\begin{lemma}[Cohomological $\delta$-functor]\label{lem:les}
$\{H^i(X,-)\}$ and $\{\Rr^i f_*\}$ are cohomological $\delta$-functors: a short
exact sequence $0\to\Ff'\to\Ff\to\Ff''\to 0$ yields a natural long exact
sequence, with $H^0=\Gamma$, $\Rr^0 f_*=f_*$. Reference: Stacks,
tag~\stacks{0156}; the $\delta$-functor API is in Mathlib
(\texttt{CategoryTheory.Abelian.RightDerived}).
\lean{AlgebraicGeometry.Scheme.higherDirectImage\_deltaFunctor}
\uses{def:sheaf-cohomology, def:higher-direct-image, lem:pf-left-exact}
\mathlibpartial
\end{lemma}

\begin{lemma}[Higher direct image as a sheafification]\label{lem:hdi-sheafification}
$\Rr^i f_*\Ff$ is the sheaf associated to the presheaf $V\mapsto
H^i(f^{-1}(V),\Ff|_{f^{-1}(V)})$ on $Y$. In particular $\Rr^i f_*\Ff|_V$
depends only on $\Ff|_{f^{-1}(V)}$, so the computation of $\Rr^i f_*$ is local
on $Y$. Reference: Stacks, Lemma 20.7.3, tag~\stacks{01E0}; Hartshorne, III.8.1.
\lean{AlgebraicGeometry.Scheme.higherDirectImage\_isSheafification}
\uses{def:sheaf-cohomology, def:higher-direct-image}
\notready
\end{lemma}

\chapter{Flasque sheaves and acyclic resolutions}\label{chap:flasque}

This chapter is the light replacement for the relative Leray spectral sequence
(Remark~\ref{rem:no-ss}); it also supplies the acyclic resolutions used to
compute cohomology in later parts.

\begin{definition}[Flasque sheaf]\label{def:flasque}
A sheaf $\Ff$ on a topological space is \emph{flasque} if $\Ff(U)\to\Ff(V)$ is
surjective for all opens $V\subseteq U$. Reference: Stacks, tag~\stacks{01AF};
Hartshorne, II.\,Ex.1.16.
\lean{TopCat.Presheaf.IsFlasque}
\notready
\end{definition}

\begin{lemma}[Injectives are flasque]\label{lem:injective-flasque}
Every injective $\Oo_X$-module is flasque. Reference: Stacks,
tag~\stacks{01EA}; Hartshorne, III.2.4.
\lean{AlgebraicGeometry.Scheme.Modules.injective\_isFlasque}
\uses{def:flasque}
\notready
\end{lemma}

\begin{lemma}[Pushforward preserves flasqueness]\label{lem:pushforward-flasque}
If $\Ff$ is flasque on $X$ and $g\colon X\to X'$ continuous, then $g_*\Ff$ is
flasque, since $(g_*\Ff)(V)=\Ff(g^{-1}V)$. Reference: Stacks, tag~\stacks{01E9}.
\lean{TopCat.Presheaf.isFlasque\_pushforward}
\uses{def:flasque}
\notready
\end{lemma}

\begin{lemma}[Flasque sheaves are acyclic]\label{lem:flasque-acyclic}
If $\Ff$ is flasque then $H^i(X,\Ff)=0$ and $\Rr^i f_*\Ff=0$ for all $i>0$.
Reference: Stacks, Lemma 20.12.3, tag~\stacks{01EB}; Hartshorne, III.2.5.
\lean{AlgebraicGeometry.Scheme.higherDirectImage\_flasque\_eq\_zero}
\uses{def:flasque, lem:les, lem:hdi-sheafification}
\notready
\end{lemma}

\begin{proof}[Proof sketch]
Flasqueness is preserved under quotients by flasque subsheaves, and if two
terms of a short exact sequence are flasque so is the third; embedding $\Ff$
in an injective (hence flasque) $\Ii$ and using the long exact sequence
(Lemma~\ref{lem:les}) gives the vanishing by dimension shifting. For
$\Rr^i f_*$ use Lemma~\ref{lem:hdi-sheafification} and the same argument on
each $f^{-1}(V)$. Faithful to Stacks, tag~\stacks{01EB}.
\end{proof}

\begin{lemma}[Acyclic resolutions compute derived functors]\label{lem:acyclic-resolution}
If $0\to\Ff\to\mathcal A^0\to\mathcal A^1\to\cdots$ is a resolution with each
$\mathcal A^k$ acyclic for $\Gamma(X,-)$ (resp. $f_*$), then
$H^i(X,\Ff)\cong H^i(\Gamma(X,\mathcal A^\bullet))$ (resp.
$\Rr^i f_*\Ff\cong H^i(f_*\mathcal A^\bullet)$). Reference: Stacks,
tag~\stacks{0156} (derived functor via acyclic resolution); Hartshorne, III.1.2A.
\lean{AlgebraicGeometry.Scheme.cohomology\_of\_acyclicResolution}
\uses{lem:flasque-acyclic}
\notready
\end{lemma}

\begin{lemma}[Acyclic pushforward lemma]\label{lem:acyclic-pushforward}
Let $X\xrightarrow{\,g\,}X'\xrightarrow{\,f\,}Y$ and let $\Bb$ be an
$\Oo_X$-module with $\Rr^q g_*\Bb=0$ for all $q>0$. Then
$\Rr^q f_*(g_*\Bb)\cong\Rr^q(f\circ g)_*\Bb$ naturally for all $q\ge 0$.
Reference: the note \cite{Note}.
\lean{AlgebraicGeometry.Scheme.higherDirectImage\_pushforward\_of\_acyclic}
\uses{lem:injective-flasque, lem:pushforward-flasque, lem:flasque-acyclic, lem:acyclic-resolution}
\notready
\end{lemma}

\begin{proof}[Proof]
Take an injective resolution $\Bb\to I^\bullet$. Each $I^k$ is flasque
(Lemma~\ref{lem:injective-flasque}); since $\Rr^q g_*\Bb=H^q(g_* I^\bullet)$,
the hypothesis says $g_* I^\bullet$ is a resolution of $g_*\Bb$, with terms
flasque (Lemma~\ref{lem:pushforward-flasque}) hence $f_*$-acyclic
(Lemma~\ref{lem:flasque-acyclic}). By Lemma~\ref{lem:acyclic-resolution} and
$f_*\circ g_*=(f\circ g)_*$,
$\Rr^q f_*(g_*\Bb)=H^q(f_*g_*I^\bullet)=H^q((f\circ g)_*I^\bullet)
=\Rr^q(f\circ g)_*\Bb$. Faithful to \cite{Note}.
\end{proof}

\chapter{Comparison: abstract \texorpdfstring{$\mathrm{rightDerived}$}{rightDerived} vs sheaf cohomology}\label{chap:bridge1}

\begin{proposition}[The abstract derived pushforward computes $\Rr^n f_*$]\label{prop:bridge-rderived}
For $f\colon X\to Y$ and a quasi-coherent $M\in\texttt{SheafOfModules}(X)$ there
is a natural isomorphism
$\bigl((\mathrm{Modules.pushforward}\,f).\mathrm{rightDerived}\,n\bigr)(M)\cong
\Rr^n f_*M$, compatible with the equivalence between geometric
$\Oo_X$-modules and \texttt{SheafOfModules}. In particular the object in
Theorem~\ref{thm:goal} is the geometric higher direct image.
\lean{AlgebraicGeometry.Scheme.Modules.rightDerivedPushforward\_iso\_higherDirectImage}
\uses{def:higher-direct-image, lem:pf-left-exact, lem:flasque-acyclic}
\notready
\end{proposition}

\begin{proof}[Proof strategy]
Both are $\delta$-functors agreeing in degree $0$ ($f_*$); it suffices to show
the abstract one is erasable/effaceable, i.e.\ computed by a class of objects
acyclic on both sides. Injective objects of \texttt{SheafOfModules} map to
flasque $\Oo_X$-modules (Lemma~\ref{lem:injective-flasque} adapted), which are
$f_*$-acyclic (Lemma~\ref{lem:flasque-acyclic}); by uniqueness of derived
functors the two $\delta$-functors are isomorphic. This is the single
comparison tying the categorical machinery of the target statement to the
geometric development.
\end{proof}

%##############################################################################
\part{\v{C}ech cohomology and vanishing on affines}\label{part:cech}
%##############################################################################

References: Stacks, Sections 20.9--20.11 and 30.2--30.4 \cite{stacks};
Hartshorne, III.3--III.4 \cite{Hartshorne}.

\chapter{The \v{C}ech complex}\label{chap:cech}

\begin{definition}[\v{C}ech complex and cohomology]\label{def:cech}
For an open cover $\mathcal U=\{U_i\}_{i\in I}$ of $X$ and an $\Oo_X$-module
$\Ff$, the \v{C}ech complex is
$\Cech^p(\mathcal U,\Ff)=\prod_{i_0<\dots<i_p}\Ff(U_{i_0\cdots i_p})$ with the
usual alternating differential, and $\check H^p(\mathcal U,\Ff)$ its
cohomology. Reference: Stacks, tag~\stacks{01ED}; Hartshorne, III.4.
\lean{AlgebraicGeometry.Scheme.cechComplex, AlgebraicGeometry.Scheme.cechCohomology}
\notready
\end{definition}

\begin{lemma}[\v{C}ech degree $0$ and functoriality]\label{lem:cech-basic}
$\check H^0(\mathcal U,\Ff)\cong\Gamma(X,\Ff)$, and a short exact sequence of
sheaves induces a long exact sequence of \v{C}ech cohomology when the cover
consists of sheafy-acyclic pieces (e.g.\ affines with quasi-coherent $\Ff$).
Reference: Stacks, tags~\stacks{01EF}, \stacks{01EO}.
\lean{AlgebraicGeometry.Scheme.cechCohomology\_zero}
\uses{def:cech}
\notready
\end{lemma}

\begin{proposition}[\v{C}ech-to-derived comparison]\label{prop:cech-to-derived}
There is a natural map $\check H^p(\mathcal U,\Ff)\to H^p(X,\Ff)$, an
isomorphism in degrees $0,1$ in general, and an isomorphism in all degrees when
every finite intersection $U_{i_0\cdots i_p}$ is $\Ff$-acyclic (Leray's
theorem / acyclic covers). Reference: Stacks, Lemma 20.11.6, tag~\stacks{01EW};
Hartshorne, III.4.5.
\lean{AlgebraicGeometry.Scheme.cechToDerived, AlgebraicGeometry.Scheme.cechToDerived\_iso\_of\_acyclicCover}
\uses{def:cech, def:sheaf-cohomology, lem:flasque-acyclic}
\notready
\end{proposition}

\chapter{Vanishing on affines and the comparison}\label{chap:affine-vanishing}

\begin{theorem}[Serre: vanishing on affines]\label{thm:affine-vanishing}
If $X=\Spec R$ is affine and $\Ff$ is quasi-coherent, then $H^i(X,\Ff)=0$ for
all $i>0$. Reference: Stacks, Lemma 30.2.2, tag~\stacks{01XB}; Hartshorne,
III.3.5 (Noetherian case) and III.3.7.
\lean{AlgebraicGeometry.cohomology\_qcoh\_affine\_eq\_zero}
\uses{def:qcoh, prop:affine-equiv, def:sheaf-cohomology, prop:cech-to-derived}
\notready
\end{theorem}

\begin{proof}[Proof outline]
By Lemma~\ref{lem:hdi-sheafification}/\ref{prop:cech-to-derived} reduce to
showing the \v{C}ech cohomology of $\widetilde M$ for the standard cover of
$\Spec R$ by basic opens $D(f_i)$ (with $\sum(f_i)=R$) vanishes above degree
$0$; the \v{C}ech complex is the Koszul-type complex computing
$\check H^p=0$ for $p>0$ because localisation is exact and the $f_i$ generate
the unit ideal. In the Noetherian case one may instead use that a quasi-coherent
sheaf on $\Spec R$ embeds in a flasque quasi-coherent sheaf and induct. Faithful
to Stacks, tag~\stacks{01XB}.
\end{proof}

\begin{corollary}[Affine covers compute cohomology]\label{cor:affine-cover-computes}
If $X$ is separated (so intersections of affine opens are affine) and
$\mathcal U$ is an affine open cover, then for quasi-coherent $\Ff$,
$\check H^p(\mathcal U,\Ff)\cong H^p(X,\Ff)$ for all $p$. Reference: Stacks,
Lemma 30.2.6, tag~\stacks{01X9}; Hartshorne, III.4.5.
\lean{AlgebraicGeometry.Scheme.cechCohomology\_iso\_cohomology\_of\_separated}
\uses{prop:cech-to-derived, thm:affine-vanishing}
\notready
\end{corollary}

\chapter{Higher direct images of quasi-coherent sheaves}\label{chap:hdi-qcoh}

\begin{proposition}[$\Rr^i f_*$ preserves quasi-coherence]\label{prop:hdi-qcoh}
If $f\colon X\to Y$ is quasi-compact and quasi-separated and $\Ff$ is
quasi-coherent, then $\Rr^i f_*\Ff$ is quasi-coherent for all $i\ge 0$.
Reference: Stacks, Lemma 30.4.5, tag~\stacks{01XJ}.
\lean{AlgebraicGeometry.Scheme.higherDirectImage\_isQuasicoherent}
\uses{def:higher-direct-image, lem:hdi-sheafification, thm:affine-vanishing, prop:qcoh-abelian}
\notready
\end{proposition}

\begin{proof}[Proof outline]
By Lemma~\ref{lem:hdi-sheafification} the question is local on $Y$, so take
$Y=\Spec A$. Cover $X$ by finitely many affines; by
Corollary~\ref{cor:affine-cover-computes} the \v{C}ech complex of quasi-coherent
sheaves computes $\Rr^i f_*\Ff$, and each \v{C}ech term is quasi-coherent
(pushforward of quasi-coherent along the affine inclusions of intersections,
using $qcqs$); quasi-coherence is closed under kernels/cokernels
(Proposition~\ref{prop:qcoh-abelian}), so the cohomology is quasi-coherent.
Faithful to Stacks, tag~\stacks{01XJ}.
\end{proof}

\begin{proposition}[Direct image over an affine base]\label{prop:hdi-affine-base}
For $A$ a ring, $f\colon X\to\Spec A$ quasi-compact quasi-separated and $\Ff$
quasi-coherent, $\Rr^i f_*\Ff\cong\widetilde{H^i(X,\Ff)}$; in particular
$\Rr^i f_*\Ff$ is coherent iff $H^i(X,\Ff)$ is a finite $A$-module. Reference:
Stacks, Lemma 30.4.6, tag~\stacks{01XK}.
\lean{AlgebraicGeometry.Scheme.higherDirectImage\_over\_affine}
\uses{prop:hdi-qcoh, prop:affine-equiv, lem:tilde-sections}
\notready
\end{proposition}

%##############################################################################
\part{Projective space and Proj}\label{part:proj}
%##############################################################################

References: Stacks, Sections 27.10--27.14 (Proj) and 30.8, 30.14--30.15
\cite{stacks}; Hartshorne, II.5, III.5 \cite{Hartshorne}; Serre \cite{Serre-FAC}.

\chapter{Twisting sheaves and graded modules on Proj}\label{chap:proj-twist}

\begin{construction}[The functor $\widetilde{(-)}$ on Proj and twists]\label{con:proj-tilde}
Let $S$ be a graded ring, finitely generated in degree $1$ over $S_0=A$, and
$X=\Proj S$ (a scheme in Mathlib). For a graded $S$-module $M$ there is a
quasi-coherent sheaf $\widetilde M$ on $X$; taking $M=S(d)$ defines the
twisting sheaves $\Oo_X(d)=\widetilde{S(d)}$, with
$\Oo_X(d)\otimes\Oo_X(e)\cong\Oo_X(d+e)$ and $\Oo_X(d)$ invertible when $S$ is
generated in degree $1$. Write $\Ff(d)=\Ff\otimes\Oo_X(d)$. Reference: Stacks,
Section 27.10, tags~\stacks{01M3}, \stacks{01MT}; Hartshorne, II.5.
\lean{AlgebraicGeometry.Proj.tildeGraded, AlgebraicGeometry.Proj.twistingSheaf}
\uses{def:qcoh, def:invertible}
\notready
\end{construction}

\begin{lemma}[Sections on the standard cover]\label{lem:proj-sections}
For $f\in S_1$, $D_+(f)=\Spec (S_{(f)})$ and
$\Oo_X(d)|_{D_+(f)}\cong\widetilde{(S(d)_{(f)})}$; the standard cover
$\{D_+(x_i)\}$ for degree-$1$ generators $x_i$ is affine, and its intersections
are the $D_+(x_{i_0}\cdots x_{i_p})$. Reference: Stacks, tag~\stacks{01M3}.
\lean{AlgebraicGeometry.Proj.twistingSheaf\_restrict}
\uses{con:proj-tilde, lem:tilde-sections}
\notready
\end{lemma}

\begin{proposition}[Every coherent sheaf on Proj comes from a graded module]\label{prop:proj-graded}
Let $A$ be Noetherian and $S$ a finitely generated graded $A$-algebra
generated in degree $1$, $X=\Proj S$. Every coherent $\Oo_X$-module is
isomorphic to $\widetilde M$ for a finitely generated graded $S$-module $M$,
and there is $d_0$ with $M_d\xrightarrow{\ \sim\ }H^0(X,\widetilde M(d))$ for
$d\ge d_0$. Reference: Stacks, Proposition 30.14.4 and Section 30.15,
tags~\stacks{01YS}, \stacks{01YW}; Hartshorne, II.5.15--5.17.
\lean{AlgebraicGeometry.Proj.coherent\_from\_gradedModule}
\uses{con:proj-tilde, def:coherent, lem:proj-sections}
\notready
\end{proposition}

\chapter{Cohomology of projective space}\label{chap:cohomology-Pn}

\begin{theorem}[Serre's computation]\label{thm:cohomology-Pn}
Let $A$ be a ring, $n\ge 0$, $\PP^n_A=\Proj A[T_0,\dots,T_n]$. Then, as graded
$A$-modules,
\[
\bigoplus_d H^0(\PP^n_A,\Oo(d))\cong A[T_0,\dots,T_n],\qquad
\bigoplus_d H^n(\PP^n_A,\Oo(d))\cong
\tfrac{1}{T_0\cdots T_n}A[\tfrac1{T_0},\dots,\tfrac1{T_n}],
\]
and $H^q(\PP^n_A,\Oo(d))=0$ for $0<q<n$; each $H^q(\PP^n_A,\Oo(d))$ is a
finite free $A$-module. Reference: Stacks, Lemma 30.8.1, tag~\stacks{01XS};
Hartshorne, III.5.1.
\lean{AlgebraicGeometry.cohomology\_projectiveSpace\_twist}
\uses{thm:affine-vanishing, cor:affine-cover-computes, lem:proj-sections}
\notready
\end{theorem}

\begin{proof}[Proof outline]
The standard cover $\{D_+(T_i)\}$ is affine with affine intersections, so by
Corollary~\ref{cor:affine-cover-computes} the cohomology is that of the
\v{C}ech complex $\Cech^\bullet(\{D_+(T_i)\},\Oo(d))$. This complex is, in each
graded piece, an explicit Koszul-type complex of localisations of the
polynomial ring; a direct computation of its cohomology gives $H^0$ (degree
$d$ polynomials), $H^n$ (the stated ``dual'' piece), and vanishing in
between. Finiteness/freeness over $A$ is read off the monomial bases. Faithful
to Stacks, tag~\stacks{01XS}.
\end{proof}

\begin{corollary}[Finiteness and Serre vanishing on $\PP^n$]\label{cor:Pn-finite-vanish}
$H^q(\PP^n_A,\Oo(d))$ is a finite $A$-module for all $q,d$, and vanishes for
$q>0$ and $d\gg 0$.
\lean{AlgebraicGeometry.cohomology\_projectiveSpace\_finite\_and\_vanishing}
\uses{thm:cohomology-Pn}
\notready
\end{corollary}

\chapter{Serre finiteness and vanishing on Proj}\label{chap:serre-proj}

\begin{lemma}[Affine finiteness, projective case]\label{lem:affine-finite-proj}
Let $A$ be Noetherian and $X=\Proj S$ projective over $\Spec A$ (with $S$ as in
Proposition~\ref{prop:proj-graded}), $\Ff$ coherent. Then $H^i(X,\Ff)$ is a
finite $A$-module for all $i$, and $H^i(X,\Ff(d))=0$ for $i>0$, $d\gg 0$;
moreover $\Ff(d)$ is globally generated for $d\gg 0$. Reference: Stacks, Lemma
30.15.2 and Section 30.15, tag~\stacks{01YW}; Hartshorne, III.5.2.
\lean{AlgebraicGeometry.Proj.cohomology\_coherent\_finite\_and\_vanishing}
\uses{thm:cohomology-Pn, prop:proj-graded, def:coherent}
\notready
\end{lemma}

\begin{proof}[Proof outline]
Write $\Ff=\widetilde M$ (Proposition~\ref{prop:proj-graded}) and present $M$
as a quotient of a finite free graded module; this yields a finite resolution
of $\Ff$ by finite direct sums of twists $\Oo_X(a_j)$ up to a coherent kernel,
and $H^\bullet$ of each such sum is finite over $A$ by
Theorem~\ref{thm:cohomology-Pn}. Descending induction on the (finite)
cohomological amplitude, using the long exact sequence
(Lemma~\ref{lem:les}), gives finiteness of $H^i(X,\Ff)$ and, after twisting by
$\Oo(d)$ with $d\gg0$, the vanishing and global generation. Faithful to Stacks,
tag~\stacks{01YW}.
\end{proof}

%##############################################################################
\part{Coherent sheaves}\label{part:coherent}
%##############################################################################

References: Stacks, Sections 30.9, 30.11, 30.12 \cite{stacks}; Hartshorne, II.5
\cite{Hartshorne}.

\chapter{Coherence on locally Noetherian schemes}\label{chap:coherent}

\begin{definition}[Coherent module]\label{def:coherent}
An $\Oo_X$-module $\Ff$ is \emph{coherent} if it is of finite type and, for
every open $U$ and every morphism $\Oo_U^{\,m}\to\Ff|_U$, the kernel is of
finite type. Reference: Stacks, tag~\stacks{01BU}.
\lean{AlgebraicGeometry.Scheme.Modules.IsCoherent}
\uses{}
\notready
\end{definition}

\begin{lemma}[Locally Noetherian source]\label{lem:X-noeth}
If $f\colon X\to S$ is of finite type and $S$ is locally Noetherian, then $X$
is locally Noetherian. In particular a proper $f$ over locally Noetherian $S$
has locally Noetherian source. Reference: Stacks, Lemma 29.15.6,
tag~\stacks{01T6}.
\lean{AlgebraicGeometry.locallyNoetherian\_of\_finiteType}
\mathlibpartial
\end{lemma}

\begin{lemma}[Coherence on locally Noetherian schemes]\label{lem:coherent-noeth}
On a locally Noetherian scheme $X$, an $\Oo_X$-module is coherent iff it is
quasi-coherent and of finite type. Reference: Stacks, Lemma 30.9.1,
tag~\stacks{01XZ}.
\lean{AlgebraicGeometry.Scheme.Modules.isCoherent\_iff\_qcoh\_finiteType}
\uses{def:coherent, def:qcoh}
\notready
\end{lemma}

\begin{lemma}[Finite presentation $=$ coherence]\label{lem:fp-iff-coh}
On a locally Noetherian scheme, an $\Oo_X$-module is of finite presentation
(in the sense of \texttt{SheafOfModules.IsFinitePresentation}) iff it is
coherent. Reference: Stacks, tag~\stacks{01XZ} together with
tag~\stacks{01BN}.
\lean{AlgebraicGeometry.Scheme.Modules.isFinitePresentation\_iff\_isCoherent}
\uses{lem:coherent-noeth, lem:qcoh-agree}
\notready
\end{lemma}

\begin{proposition}[Coherent sheaves form a Serre--abelian subcategory]\label{prop:coherent-abelian}
On a locally Noetherian scheme, kernels, cokernels, images and extensions of
coherent sheaves are coherent; in a short exact sequence, if two terms are
coherent so is the third. Reference: Stacks, Lemmas 30.9.2, 30.9.3,
tags~\stacks{01Y0}, \stacks{01Y1}.
\lean{AlgebraicGeometry.Scheme.Modules.isCoherent\_of\_ses}
\uses{def:coherent, lem:coherent-noeth, prop:qcoh-abelian}
\notready
\end{proposition}

\chapter{Closed immersions and support}\label{chap:closed-support}

\begin{lemma}[Pushforward along a closed immersion]\label{lem:closed-immersion-acyclic}
Let $i\colon Z\to X$ be a closed immersion of locally Noetherian schemes and
$\Gg$ a coherent $\Oo_Z$-module. Then $i_*\Gg$ is coherent, $i_*$ is exact, and
$\Rr^p i_*\Gg=0$ for all $p>0$. Reference: Stacks, Lemma 30.9.9,
tag~\stacks{01Y6}; Liu, 5.1.14(d) \cite{Liu}.
\lean{AlgebraicGeometry.Scheme.closedImmersion\_pushforward\_coherent\_and\_acyclic}
\uses{def:coherent, def:higher-direct-image, lem:flasque-acyclic}
\notready
\end{lemma}

\begin{proof}[Proof sketch]
Coherence of $i_*\Gg$ is local and reduces to the affine statement that a
finite module over $R/I$ is finite over $R$. Exactness and the vanishing of
$\Rr^p i_*$ for $p>0$ hold because $i$ is a homeomorphism onto a closed subset,
so $i_*$ is exact on the underlying sheaves and $i_*$ of a flasque sheaf is
flasque, hence acyclic (Lemma~\ref{lem:flasque-acyclic}). Faithful to Stacks,
tag~\stacks{01Y6}.
\end{proof}

\begin{lemma}[Support and generic stalks]\label{lem:support}
For a coherent $\Ff$ on a Noetherian scheme, $\Supp(\Ff)$ is closed; if
$Z\subseteq X$ is integral closed with generic point $\xi$ then
$\Oo_{X,\xi}=\kk(\xi)$ is a field, so any coherent $\Gg$ supported on $Z$ has
$\Gg_\xi$ a finite-dimensional $\kk(\xi)$-vector space, and $\Gg_\xi$ is
automatically annihilated by $\mathfrak m_\xi=0$. Reference: Stacks,
tags~\stacks{01B4}, \stacks{01YB}.
\lean{AlgebraicGeometry.Scheme.Modules.support\_closed, AlgebraicGeometry.Scheme.Modules.genericStalk\_finrank}
\uses{def:coherent}
\notready
\end{lemma}

%##############################################################################
\part{Projective morphisms}\label{part:projmor}
%##############################################################################

References: Stacks, Sections 29.40 (ample), 30.16 \cite{stacks}; Hartshorne,
II.7, III.8.8 \cite{Hartshorne}.

\chapter{Ample invertible sheaves and projective morphisms}\label{chap:ample}

\begin{definition}[Relatively ample; projective morphism]\label{def:projective-morphism}
An invertible $\Oo_X$-module $\Ll$ is \emph{$f$-relatively ample} (for
$f\colon X\to S$) if, locally on $S=\bigcup\Spec A_i$, $X_i\to\Spec A_i$ is
isomorphic over $A_i$ to a locally closed immersion into some $\PP^{n_i}_{A_i}$
pulling $\Oo(1)$ back to $\Ll$. A morphism $f\colon X\to S$ is
\emph{projective} if it factors as a closed immersion $X\hookrightarrow\PP^n_S$
followed by the projection (equivalently, admits an $f$-ample $\Ll$ with $f$
proper). Reference: Stacks, Section 29.40, tag~\stacks{01VG}, and
tag~\stacks{01W8}; Hartshorne, II.7.
\lean{AlgebraicGeometry.Scheme.IsRelativelyAmple, AlgebraicGeometry.IsProjective}
\uses{def:invertible}
\notready
\end{definition}

\begin{lemma}[Pullback of $\Oo(1)$ is relatively ample]\label{lem:pullback-ample}
If $i\colon X\to\PP^n_S$ is a closed immersion over $S$ then
$i^*\Oo_{\PP^n_S}(1)$ is $f$-relatively ample, where $f\colon X\to S$;
moreover $f$ is projective and proper. Reference: Stacks, Lemma 29.40.7,
tag~\stacks{02NR}.
\lean{AlgebraicGeometry.Scheme.relativelyAmple\_pullback\_twist}
\uses{def:projective-morphism}
\notready
\end{lemma}

\chapter{Higher direct images along projective morphisms}\label{chap:proj-hdi}

\begin{lemma}[Graded finiteness along a projective morphism]\label{lem:graded-finiteness}
Let $S$ be Noetherian affine, $f\colon X\to\Spec A$ projective with
$f$-relatively ample $\Ll$, and $\Ff$ coherent. Then
$\bigoplus_d H^p(X,\Ff\otimes\Ll^{\otimes d})$ is a finite graded module over
the section ring, for each $p$; in particular each $H^p(X,\Ff)$ is a finite
$A$-module and $H^p(X,\Ff\otimes\Ll^{\otimes d})=0$ for $p>0$, $d\gg 0$.
Reference: Stacks, Lemma 30.16.1 and Lemma 30.19.3, tags~\stacks{02O0},
\stacks{0897}.
\lean{AlgebraicGeometry.Scheme.gradedCohomology\_finite\_projective}
\uses{lem:affine-finite-proj, prop:proj-graded, def:projective-morphism}
\notready
\end{lemma}

\begin{lemma}[Vanishing after a large relatively ample twist]\label{lem:proj-hdi-vanish}
Let $S$ be Noetherian, $f\colon X\to S$ proper, $\Ll$ an $f$-relatively ample
invertible sheaf, $\Ff$ coherent. Then there is $n_0$ with
$\Rr^p f_*(\Ff\otimes\Ll^{\otimes n})=0$ for all $p>0$ and $n\ge n_0$.
Reference: Stacks, Lemma 30.16.2, tag~\stacks{02O1}.
\lean{AlgebraicGeometry.Scheme.higherDirectImage\_vanishing\_relativelyAmple}
\uses{lem:graded-finiteness, prop:hdi-affine-base, lem:hdi-sheafification}
\notready
\end{lemma}

\begin{proof}[Proof sketch]
The statement is local on $S$ (Lemma~\ref{lem:hdi-sheafification}), so take
$S=\Spec A$ Noetherian. By Proposition~\ref{prop:hdi-affine-base},
$\Rr^p f_*(\Ff\otimes\Ll^{\otimes n})=\widetilde{H^p(X,\Ff\otimes
\Ll^{\otimes n})}$, which vanishes for $p>0$, $n\gg0$ by the Serre vanishing in
Lemma~\ref{lem:graded-finiteness}. Faithful to Stacks, tag~\stacks{02O1}.
\end{proof}

\begin{lemma}[Coherence in the projective case]\label{lem:proj-hdi-coherent}
Let $S$ be locally Noetherian, $f\colon X\to S$ projective, $\Ff$ coherent.
Then $\Rr^p f_*\Ff$ is coherent for all $p\ge 0$. Reference: Stacks, Lemma
30.16.3, tag~\stacks{02O4}; Hartshorne, III.8.8.
\lean{AlgebraicGeometry.Scheme.higherDirectImage\_coherent\_of\_projective}
\uses{lem:graded-finiteness, prop:hdi-affine-base, lem:coherent-noeth}
\notready
\end{lemma}

\begin{proof}[Proof sketch]
Local on $S$, reduce to $S=\Spec A$ Noetherian; then
$\Rr^p f_*\Ff=\widetilde{H^p(X,\Ff)}$
(Proposition~\ref{prop:hdi-affine-base}), and $H^p(X,\Ff)$ is a finite
$A$-module by Lemma~\ref{lem:graded-finiteness}/\ref{lem:affine-finite-proj},
so $\Rr^p f_*\Ff$ is coherent (Lemma~\ref{lem:coherent-noeth}). Faithful to
Stacks, tag~\stacks{02O4}.
\end{proof}

%##############################################################################
\part{D\'evissage}\label{part:devissage}
%##############################################################################

References: Stacks, Section 30.12 \cite{stacks}.

\chapter{D\'evissage of coherent sheaves}\label{chap:devissage}

\begin{lemma}[Generation by pushforwards from integral subschemes]\label{lem:devissage-gen}
Let $X$ be Noetherian and $\mathcal P$ a property of coherent sheaves that
holds for $(Z\hookrightarrow X)_*\Ii$ for every integral closed
$Z\subseteq X$ and every coherent ideal sheaf $\Ii\subseteq\Oo_Z$. Then
$\mathcal P$ holds for all coherent sheaves. Reference: Stacks, Lemmas 30.12.4,
30.12.5, tags~\stacks{01YG}, \stacks{01YH}.
\lean{AlgebraicGeometry.Scheme.devissage\_generation}
\uses{def:coherent, prop:coherent-abelian}
\notready
\end{lemma}

\begin{theorem}[D\'evissage]\label{thm:devissage}
Let $X$ be Noetherian and $\mathcal P$ a property of coherent
$\Oo_X$-modules such that (1) in any short exact sequence, if two terms satisfy
$\mathcal P$ so does the third; and (2) for every integral closed
$Z\subseteq X$ with generic point $\xi$ there is a coherent $\Gg$ with
$\Supp\Gg=Z$, $\dim_{\kk(\xi)}\Gg_\xi=1$, and $\mathcal P(\Gg)$. Then
$\mathcal P$ holds for every coherent $\Oo_X$-module. Reference: Stacks, Lemma
30.12.6, tag~\stacks{01YI}.
\lean{AlgebraicGeometry.Scheme.coherent\_devissage}
\uses{lem:devissage-gen, prop:coherent-abelian, lem:support}
\notready
\end{theorem}

\begin{proof}[Proof sketch]
By Lemma~\ref{lem:devissage-gen} it suffices to treat
$(Z\hookrightarrow X)_*\Ii$. If $\mathcal P$ failed for some such sheaf,
Noetherian induction on the support gives a minimal integral closed $Z_0$ on
which it fails, while it holds on every proper closed subscheme of $Z_0$;
hypothesis (2) for $Z_0$ and the 2-out-of-3 property (1) applied to a suitable
short exact sequence relating $\Ii$, $\Gg$ and sheaves supported on lower-dimensional
closed subsets yield a contradiction. Faithful to Stacks, tag~\stacks{01YI}.
\end{proof}

%##############################################################################
\part{Chow's lemma}\label{part:chow}
%##############################################################################

References: Stacks, Section 30.18 \cite{stacks}; EGA II, 5.6.1 \cite{EGA2}.

\chapter{Chow's lemma}\label{chap:chow}

\begin{theorem}[Chow's lemma]\label{thm:chow}
Let $S$ be Noetherian and $f\colon X\to S$ separated of finite type. Then there
exist $n$, a proper surjective $\pi\colon X'\to X$ that is an isomorphism over a
dense open $U\subseteq X$, and an immersion $X'\hookrightarrow\PP^n_S$ over $S$.
If $f$ is proper then $X'\to S$ is projective. Reference: Stacks, Lemma
30.18.1, tag~\stacks{0200}.
\lean{AlgebraicGeometry.chowLemma}
\uses{def:projective-morphism}
\notready
\end{theorem}

\begin{proof}[Proof outline]
Cover $X$ by finitely many affine opens $U_i$ containing all generic points,
so $U=\bigcap U_i$ is dense; embed each $U_i$ as a scheme-theoretically dense
open of a closed subscheme $Z_i\subseteq\PP^{n_i}_S$. Let $Z$ be the
scheme-theoretic image of $U\to\prod_S Z_i$; the projections give proper
$p_i\colon Z\to Z_i$, and gluing the $p_i^{-1}(U_i)$ produces
$\pi\colon X'\to X$ proper and surjective, an isomorphism over $U$, with $X'$
placed in one $\PP^n_S$ by the Segre embedding. Faithful to Stacks,
tag~\stacks{0200}.
\end{proof}

\begin{remark}[The graph is a closed immersion; ampleness]\label{rem:chow-graph}
If $i\colon X'\to\PP^m_S$ is the immersion and $\pi\colon X'\to X$ proper, then
$(i,\pi)\colon X'\to\PP^m_X$ is a closed immersion, so
$\Ll:=i^*\Oo(1)\cong(i,\pi)^*\Oo_{\PP^m_X}(1)$ is $\pi$-relatively ample; and
$X'\to S$ is projective. Reference: Stacks, Remark 30.18.2, tag~\stacks{0201};
Lemma~\ref{lem:pullback-ample}.
\lean{AlgebraicGeometry.chowLemma\_graph\_closedImmersion}
\uses{thm:chow, lem:pullback-ample}
\notready
\end{remark}

%##############################################################################
\part{The finiteness theorem}\label{part:finiteness}
%##############################################################################

\chapter{Assembly}\label{chap:finiteness}

\begin{theorem}[Coherence of higher direct images]\label{thm:finiteness}
Let $S$ be locally Noetherian, $f\colon X\to S$ proper, $\Ff$ coherent. Then
$\Rr^p f_*\Ff$ is coherent for all $p\ge 0$. Reference: Stacks, Proposition
30.19.1, tag~\stacks{02O5}; EGA III, 3.2.1.
\lean{AlgebraicGeometry.Scheme.higherDirectImage\_coherent\_of\_isProper}
\uses{thm:devissage, thm:chow, rem:chow-graph, lem:proj-hdi-coherent, lem:proj-hdi-vanish, lem:acyclic-pushforward, lem:closed-immersion-acyclic, prop:coherent-abelian, prop:hdi-qcoh, lem:X-noeth, lem:support}
\notready
\end{theorem}

\begin{proof}
Coherence of $\Rr^p f_*\Ff$ is local on $S$ (Lemma~\ref{lem:hdi-sheafification}
and Proposition~\ref{prop:hdi-qcoh}), so assume $S$ Noetherian; then $X$ is
Noetherian (Lemma~\ref{lem:X-noeth}). Put
$\mathcal P(\Ff)\Leftrightarrow$ ``$\Rr^p f_*\Ff$ coherent for all $p$''.

\emph{Hypothesis (1) of Theorem~\ref{thm:devissage}.} For a short exact
sequence of coherent sheaves, the long exact sequence of the $\Rr^p f_*$
(Lemma~\ref{lem:les}) places each $\Rr^p f_*$ of the third sheaf between two
coherent terms; by Proposition~\ref{prop:coherent-abelian} (2-out-of-3) it is
coherent. So $\mathcal P$ satisfies 2-out-of-3.

\emph{Hypothesis (2).} Fix an integral closed $Z\subseteq X$ with generic
point $\xi$, closed immersion $j\colon Z\to X$, and $g=f\circ j\colon Z\to S$.
It suffices to produce coherent $\Gg$ on $Z$ with (a) $\dim_{\kk(\xi)}\Gg_\xi
=1$ and (b) $\Rr^p g_*\Gg$ coherent for all $p$: then $\Ff_0:=j_*\Gg$ is
coherent (Lemma~\ref{lem:closed-immersion-acyclic}), and since
$\Rr^q j_*\Gg=0$ for $q>0$ the acyclic pushforward
Lemma~\ref{lem:acyclic-pushforward} gives
$\Rr^p f_*\Ff_0\cong\Rr^p(f\circ j)_*\Gg=\Rr^p g_*\Gg$, coherent by (b), so
$\mathcal P(\Ff_0)$; and $\Supp\Ff_0=Z$, $(\Ff_0)_\xi=\Gg_\xi$ is
$1$-dimensional, matching Theorem~\ref{thm:devissage}(2).

Construct $\Gg$ by Chow's lemma (Theorem~\ref{thm:chow}) applied to
$g\colon Z\to S$: a proper surjective $\pi\colon Z'\to Z$, isomorphism over a
dense open $U\subseteq Z$, with $g'=g\circ\pi\colon Z'\to S$ projective and an
invertible $\Ll$ on $Z'$ that is $\pi$-relatively ample
(Remark~\ref{rem:chow-graph}). By Lemma~\ref{lem:proj-hdi-vanish} choose
$n\ge1$ with $\Rr^q\pi_*\Ll^{\otimes n}=0$ for $q>0$, and set
$\Gg=\pi_*\Ll^{\otimes n}$. Then:
\begin{itemize}
\item $\Gg$ is coherent, being $\pi_*\Ll^{\otimes n}=\Rr^0\pi_*\Ll^{\otimes n}$
for the projective $\pi$ (Lemma~\ref{lem:proj-hdi-coherent});
\item (a) holds: over $U$, $\pi$ is an isomorphism so
$\Gg|_U=\Ll^{\otimes n}|_U$ is invertible, giving $\Supp\Gg=Z$ and
$\dim_{\kk(\xi)}\Gg_\xi=1$ (Lemma~\ref{lem:support});
\item (b) holds by Lemma~\ref{lem:acyclic-pushforward} applied to
$Z'\xrightarrow{\pi}Z\xrightarrow{g}S$ with $\Bb=\Ll^{\otimes n}$: as
$\Rr^q\pi_*\Ll^{\otimes n}=0$ for $q>0$,
$\Rr^p g_*\Gg\cong\Rr^p g'_*\Ll^{\otimes n}$, coherent for all $p$ because
$g'$ is projective (Lemma~\ref{lem:proj-hdi-coherent}).
\end{itemize}

By Theorem~\ref{thm:devissage}, $\mathcal P(\Ff)$ holds for all coherent
$\Ff$. Only the single vanishing $\Rr^q\pi_*\Ll^{\otimes n}=0$ ($q>0$) is
used, coherence being transported through
Lemma~\ref{lem:acyclic-pushforward}; no spectral sequence appears
(Remark~\ref{rem:no-ss}). Faithful to Stacks, tag~\stacks{02O5}.
\end{proof}

%##############################################################################
\part{The target Lean statement}\label{part:bridge}
%##############################################################################

\chapter{From coherence to finite presentation}\label{chap:conclude}

\begin{theorem}[Target]\label{thm:target-final}
Let $f\colon X\to Y$ be proper, $Y$ locally Noetherian, $M$ an $\Oo_X$-module
of finite presentation. Then for all $n$,
$\bigl((\mathrm{Modules.pushforward}\,f).\mathrm{rightDerived}\,n\bigr)(M)$ is
of finite presentation. This is Theorem~\ref{thm:goal}.
\lean{AlgebraicGeometry.Scheme.rightDerivedPushforward\_isFinitePresentation\_of\_isProper}
\uses{thm:finiteness, prop:bridge-rderived, lem:fp-iff-coh, lem:X-noeth}
\notready
\end{theorem}

\begin{proof}
$f$ proper $\Rightarrow$ of finite type, so $X$ is locally Noetherian
(Lemma~\ref{lem:X-noeth}); by Lemma~\ref{lem:fp-iff-coh}, $M$ is coherent. By
Proposition~\ref{prop:bridge-rderived} the derived object is $\Rr^n f_*M$,
coherent by Theorem~\ref{thm:finiteness}; applying
Lemma~\ref{lem:fp-iff-coh} on the locally Noetherian base $Y$ returns finite
presentation.
\end{proof}

%##############################################################################
\part{Cohomology and base change}\label{part:basechange}
%##############################################################################

\chapter{Cohomology and base change}\label{chap:basechange}

References: D.~Mumford, \emph{Abelian Varieties}, \S II.5, p.~53; EGA III,
7.7--7.9; Hartshorne, III.12; Stacks, tag~\stacks{0A1H}. This companion result
(quoted as Theorem~0.1 of Buzzard's \emph{Explicit models for modular curves})
is built on the same cohomological substrate as the coherence theorem
(Part~\ref{part:cohomology}), together with relative flatness and fibrewise
cohomology.

\begin{definition}[Relative flatness and fibrewise cohomology]\label{def:flat-fiber}
For $f\colon X\to Y$ and an $\Oo_X$-module $\Ff$, say $\Ff$ is \emph{flat over
$Y$} if each stalk of $\Ff$ is flat over the corresponding local ring of $Y$,
and that \emph{$H^n(X_y,\Ff_y)=0$} if the cohomology of the restriction of $\Ff$
to the fibre $X_y$ vanishes in degree $n$, for every $y\in Y$.
\lean{AlgebraicGeometry.Scheme.Modules.FlatOver, AlgebraicGeometry.Scheme.Modules.FiberCohomologyVanishes}
\uses{def:coherent, def:higher-direct-image}
\notready
\end{definition}

\begin{theorem}[Cohomology and base change; Mumford, \emph{Ab.\ Var.}\ p.~53]\label{thm:base-change}
Let $f\colon X\to Y$ be a proper morphism, $Y$ affine, and $\Ff$ a coherent
$\Oo_X$-module that is flat over $Y$. If, for some $n$, $H^n(X_y,\Ff_y)=0$ for
all $y\in Y$, then $\Rr^{n-1}f_*\Ff$ commutes with all base changes.
\lean{AlgebraicGeometry.Scheme.higherDirectImage\_commutesWithBaseChange\_of\_fiberVanishing}
\uses{def:flat-fiber, def:higher-direct-image, thm:finiteness, lem:hdi-sheafification}
\notready
\end{theorem}

\begin{remark}
The proof (Mumford, \emph{loc.\ cit.}) studies the complex of $A$-modules
computing $\Rr^\bullet f_*\Ff$ over $Y=\Spec A$ and applies the exchange/base-change
criterion: fibrewise vanishing in degree $n$ forces the relevant boundary map to
be surjective, whence $\Rr^{n-1}f_*$ is right exact and its formation commutes
with base change. This uses the finiteness Theorem~\ref{thm:finiteness} to keep
the complex in coherent modules.
\end{remark}

%##############################################################################
\part{Dependency graph, critical path, and work packages}\label{part:summary}
%##############################################################################

\chapter{Critical path and packaging}

\section{Critical path to Theorem~\ref{thm:finiteness}}
\[
\begin{aligned}
&\text{Ch.\ref{chap:tilde} tilde} \Rightarrow \text{Ch.\ref{chap:qcoh} QCoh}
\Rightarrow \text{Ch.\ref{chap:cohomology-def} cohomology} \\
&\Rightarrow \text{Ch.\ref{chap:cech}--\ref{chap:affine-vanishing} \v{C}ech + affine vanishing}
\Rightarrow \text{Ch.\ref{chap:hdi-qcoh} qcoh of }\Rr^i f_* \\
&\Rightarrow \text{Ch.\ref{chap:cohomology-Pn} coh of }\PP^n
\Rightarrow \text{Ch.\ref{chap:proj-hdi} projective case} \\
&\text{and independently Ch.\ref{chap:flasque} flasque, Ch.\ref{chap:devissage} d\'evissage, Ch.\ref{chap:chow} Chow} \\
&\Rightarrow \text{Ch.\ref{chap:finiteness} finiteness} \Rightarrow \text{Ch.\ref{chap:conclude} target.}
\end{aligned}
\]

\section{Suggested work packages (parallelisable)}
\begin{description}[leftmargin=1.4em]
\item[WP1 Affine QCoh.] Ch.\ref{chap:tilde}--\ref{chap:qcoh}: complete the
$\widetilde{(-)}$ API in \texttt{Modules.Tilde}, the affine equivalence, and
the geometric quasi-coherence predicate. Prerequisite for everything.
\item[WP2 Cohomology core.] Ch.\ref{chap:cohomology-def}--\ref{chap:flasque}:
$H^i$, $\Rr^i f_*$, $\delta$-functor structure, sheafification description,
and the flasque package. Independent of WP1 except
Lemma~\ref{lem:hdi-sheafification}.
\item[WP3 \v{C}ech + vanishing.] Ch.\ref{chap:cech}--\ref{chap:hdi-qcoh}:
needs WP1 and WP2; yields Serre vanishing, affine-cover computation, and
quasi-coherence of $\Rr^i f_*$.
\item[WP4 Projective space.] Ch.\ref{chap:proj-twist}--\ref{chap:serre-proj}:
needs WP3; the twisting sheaves and Theorem~\ref{thm:cohomology-Pn}.
\item[WP5 Coherence.] Ch.\ref{chap:coherent}--\ref{chap:closed-support}:
needs WP1; the coherence predicate, finite-presentation equivalence, closed
immersions, support.
\item[WP6 Projective morphisms.] Ch.\ref{chap:ample}--\ref{chap:proj-hdi}:
needs WP4, WP5.
\item[WP7 D\'evissage.] Ch.\ref{chap:devissage}: needs WP5. Parallel to WP6.
\item[WP8 Chow.] Ch.\ref{chap:chow}: needs scheme-theoretic image and Segre;
largely independent morphism-theory package.
\item[WP9 Assembly + bridge.] Ch.\ref{chap:finiteness}--\ref{chap:conclude}:
needs WP2 (flasque/acyclic pushforward and Prop.\ref{prop:bridge-rderived}),
WP6, WP7, WP8.
\end{description}

\section{Largest single items}
By anticipated effort: Theorem~\ref{thm:affine-vanishing} (Serre vanishing on
affines) and the \v{C}ech comparison (WP3); Theorem~\ref{thm:cohomology-Pn}
(explicit cohomology of $\PP^n$, WP4); Theorem~\ref{thm:chow} (Chow's lemma,
WP8); and Proposition~\ref{prop:bridge-rderived} (the
\texttt{rightDerived}/sheaf-cohomology comparison, WP2). None requires a
spectral sequence.

\section{Items closest to existing Mathlib}
Lemma~\ref{lem:X-noeth}, Lemma~\ref{lem:pf-left-exact}, Lemma~\ref{lem:les}
(via the \texttt{RightDerived} $\delta$-functor API), and the tilde
sections/stalks (Lemma~\ref{lem:tilde-sections}) are the shortest to land,
being extensions of declarations already present
(Chapter~\ref{chap:baseline}).

\begin{thebibliography}{9}
\bibitem{EGA2} A.~Grothendieck, J.~Dieudonn\'e, EGA II, Publ.\ Math.\ IH\'ES 8
(1961), Th\'eor\`eme 5.6.1 (Chow's lemma).
\bibitem{EGA3} A.~Grothendieck, J.~Dieudonn\'e, EGA III, Publ.\ Math.\ IH\'ES 11
(1961), Th\'eor\`eme 3.2.1.
\bibitem{Hartshorne} R.~Hartshorne, Algebraic Geometry, GTM 52, Springer, 1977,
Chapters II.5--7, III.1--5, III.8.
\bibitem{Vakil} R.~Vakil, The Rising Sea: Foundations of Algebraic Geometry,
\S18.9.
\bibitem{Serre-FAC} J.-P.~Serre, Faisceaux alg\'ebriques coh\'erents, Ann.\ of
Math.\ 61 (1955), 197--278.
\bibitem{Liu} Q.~Liu, Algebraic Geometry and Arithmetic Curves, Oxford, 2002,
5.1.14.
\bibitem{Note} Note on avoiding spectral sequences in EGA III.3.2.1 (attached).
\bibitem{stacks} The Stacks Project Authors, The Stacks Project,
\url{https://stacks.math.columbia.edu}: Ch.~26 (tag~\stacks{01HR}), Ch.~27
(tag~\stacks{01M3}), Ch.~20 (tag~\stacks{01DW}), Ch.~30 (tag~\stacks{01X6});
Prop.~30.19.1 (tag~\stacks{02O5}).
\end{thebibliography}

\end{document}
